%% ============================================================
%%  CHAPTER 3 — MEDICAL AND PHYSIOLOGICAL FOUNDATIONS
%% ============================================================
\chapter{Medical and Physiological Foundations}
\label{ch:medical}

This chapter establishes the medical and physiological basis for the
artifact-based control paradigm adopted in this project. Whereas the
preceding chapters introduced the problem and surveyed prior work,
this chapter explains \emph{why} the chosen voluntary artifacts ---
eye blinks, jaw clinching, and lateral bruxism --- produce reliable,
distinguishable electrical signals at the scalp, and why this paradigm
is particularly well suited to the target population of upper-limb
amputees with peripheral nerve damage.

Section~\ref{sec:med_signals} distinguishes the neural and muscular
signals captured by scalp electrodes.
Section~\ref{sec:med_blink} describes the neurophysiology of the eye
blink and its electrical signature.
Section~\ref{sec:med_jaw} describes the neuromuscular basis of jaw
clinching and bruxism.
Section~\ref{sec:med_lateral} explains the anatomical basis of the
left/right bruxism asymmetry.
Section~\ref{sec:med_intact} argues why these pathways remain intact
in the target patient population.
Section~\ref{sec:med_safety} discusses the clinical safety and
comfort considerations of repeated voluntary artifact production.
Section~\ref{sec:med_rationale} synthesises the clinical rationale
for choosing an artifact-based interface over motor imagery for this
population. Section~\ref{sec:med_summary} summarises the chapter.

\medskip
\noindent\textit{Note: The physiological and clinical content of this
chapter was developed in consultation with a collaborating
specialist in neurology, who reviewed the anatomical and
neurophysiological claims for accuracy.}

%% ----------------------------------------------------------
\section{Neural versus Muscular Signals at the Scalp}
\label{sec:med_signals}
%% ----------------------------------------------------------

Electroencephalography records the summed electrical activity of the
brain through electrodes placed on the scalp. However, the scalp
potentials measured by an \acrshort{eeg} device are not exclusively
of cortical origin. They are a superposition of several distinct
biological sources, each with characteristic amplitude, frequency
content, and spatial distribution:

\begin{itemize}[leftmargin=*]
    \item \textbf{Cortical activity (true EEG):} The synchronised
    post-synaptic potentials of large populations of cortical
    neurons, measured in microvolts (typically 1--\SI{100}{\micro\volt}).
    This is the signal of interest in conventional motor-imagery
    \acrshort{bci}.

    \item \textbf{Ocular potentials (\acrshort{eog}):} Electrical changes
    produced by movement of the eyeball and eyelid, dominated by the
    corneoretinal potential. These appear most strongly at frontal
    electrodes and can reach tens to hundreds of microvolts.

    \item \textbf{Muscle activity (EMG):} The electrical activity of
    contracting skeletal muscle, particularly the facial, jaw, and
    neck muscles near the electrode sites. Surface \acrshort{emg}
    is broadband (extending well above \SI{30}{\hertz}) and can be
    an order of magnitude larger than cortical \acrshort{eeg}.
\end{itemize}

In conventional \acrshort{eeg} analysis, ocular and muscular signals
are regarded as \emph{artifacts} --- unwanted contamination to be
removed before the cortical signal can be studied~\cite{croft2000removal,
fatourechi2007emg}. The central insight of the present project is to
\emph{invert} this relationship: because these artifact signals are
large, reproducible, and under direct voluntary control, they are
deliberately produced by the user as command signals. From a control
standpoint, what matters is not whether a signal is cortical or
muscular in origin, but whether the user can generate it reliably and
the system can detect it robustly. The voluntary artifacts used here
satisfy both criteria far more readily than motor-imagery rhythms.

%% ----------------------------------------------------------
\section{Neurophysiology of the Eye Blink}
\label{sec:med_blink}
%% ----------------------------------------------------------

\subsection{Anatomy and Innervation}

The eye blink is produced principally by the \textit{orbicularis
oculi}, a sphincter-like muscle encircling the eye that closes the
eyelid when it contracts. The orbicularis oculi is innervated by the
temporal and zygomatic branches of the \textbf{facial nerve} (cranial
nerve VII). A voluntary blink is initiated by motor cortex via the
corticobulbar tract, which projects to the facial motor nucleus in
the pons; reflexive blinking additionally involves brainstem circuits.
Because the command originates supranuclearly and the effector is a
cranial muscle, the blink pathway is entirely independent of the
spinal cord and the peripheral nerves of the upper limb.

\subsection{The Corneoretinal Potential}

The dominant electrical contribution of a blink is not muscular but
ocular. The eyeball behaves as a fixed electrical dipole: the cornea
is electropositive relative to the retina, a standing voltage known
as the \textbf{corneoretinal potential}. As the eyelid sweeps down
over the cornea during a blink, it alters the conductive path between
this dipole and the frontal scalp electrodes, producing a large,
characteristic deflection at Fp1 and Fp2. A secondary contribution
comes from a brief burst of orbicularis oculi \acrshort{emg} during
lid closure. The combined event produces a deflection of approximately
50--\SI{300}{\micro\volt} at the frontal poles --- one to two orders
of magnitude larger than motor-imagery modulations --- lasting
roughly 150--\SI{400}{\milli\second}~\cite{croft2000removal,
usakli2010mouse}.

\subsection{Why Blink Counting Works --- and Its Limits}

Because each individual blink produces one discrete, time-localised
deflection at Fp1/Fp2, a sequence of blinks produces a corresponding
sequence of peaks. A single blink yields one peak; a double blink, two
peaks in rapid succession. This makes the \emph{number of voluntary
blinks} a naturally discrete control variable --- the physiological
basis for using single and double blinks as two separate commands.

In practice, however, the boundary between one blink and two is not
always clean: a subject intending a single blink may inadvertently
produce a rapid pair, and the inter-blink interval of a deliberate
double blink varies from instance to instance. Simple peak-counting is
therefore fragile. The more robust signal is the overall
\emph{morphology} of the blink waveform --- its shape over the full
epoch --- rather than a discrete peak count. This physiological
observation directly motivates the shape-based Dynamic Time Warping
approach to blink classification developed in Chapter~9, which compares
whole-waveform shape rather than counting events.

%% ----------------------------------------------------------
\section{Neuromuscular Basis of Jaw Clinching and Bruxism}
\label{sec:med_jaw}
%% ----------------------------------------------------------

\subsection{Muscles of Mastication}

Jaw movements are driven by the muscles of mastication, principally
the \textbf{masseter} and the \textbf{temporalis}. The masseter is a
thick muscle spanning the cheek from the zygomatic arch to the angle
of the mandible; the temporalis is a broad, fan-shaped muscle covering
the temporal region of the skull, lying directly beneath the T3 (left)
and T4 (right) electrode positions of the 10--20 system. Both muscles
are innervated by the mandibular division (V3) of the \textbf{trigeminal
nerve} (cranial nerve V) --- again, a cranial pathway entirely
independent of the spinal cord and upper-limb nerves.

\subsection{Jaw Clinching as a Bilateral EMG Source}

A jaw clinch --- a sustained, forceful bite --- produces a strong,
sustained contraction of the masseter and temporalis muscles on
\emph{both} sides of the head simultaneously. The resulting surface
\acrshort{emg} appears as a broadband, high-amplitude signal at the
temporal electrodes T3 and T4. Because the clinch is bilaterally
symmetric, the amplitudes at T3 and T4 are approximately equal. This
sustained, high-energy, bilateral signature makes the clinch the most
reliably detected of all five gestures, which is why it is mapped to
the safety-critical ``close hand'' command (Chapter~4, hardware).

\subsection{Bruxism as a Lateralisable Gesture}

Bruxism refers to grinding or lateral movement of the jaw. Unlike the
symmetric clinch, a deliberate lateral grind shifts the mandible to
one side and preferentially recruits the temporalis on that side. This
lateral asymmetry is the physiological key that allows left and right
bruxism to be distinguished, as described in the next section.

%% ----------------------------------------------------------
\section{Anatomical Basis of Left/Right Lateral Asymmetry}
\label{sec:med_lateral}
%% ----------------------------------------------------------

The ability to discriminate left from right bruxism rests on a
straightforward anatomical fact: the temporalis muscles on the two
sides of the head lie beneath different electrodes (T3 on the left,
T4 on the right), and a lateral grinding motion does not recruit them
equally.

When the subject grinds the jaw to the \textbf{left}, the left
temporalis is preferentially activated, producing a larger
\acrshort{emg} amplitude at T3 than at T4. When the subject grinds to
the \textbf{right}, the right temporalis dominates, producing a larger
amplitude at T4 than at T3. A symmetric clinch, by contrast, activates
both roughly equally.

This three-way pattern --- left-dominant, right-dominant, and balanced
--- maps cleanly onto a single scalar quantity: the normalised
amplitude difference between T4 and T3. This is precisely the T3--T4
asymmetry index introduced as a feature in Chapter~8, and it is the
direct algorithmic expression of the underlying muscular anatomy. The
physiological reality (which temporalis is working harder) is captured
by a single number, which in turn separates the two bruxism commands.

%% ----------------------------------------------------------
\section{Preservation of Pathways in the Target Population}
\label{sec:med_intact}
%% ----------------------------------------------------------

The clinical motivation for this project is to serve upper-limb
amputees, particularly those with extensive peripheral nerve damage
for whom conventional myoelectric prostheses are ineffective. The
viability of the artifact-based approach for this population rests on
a crucial anatomical observation: \textbf{the cranial pathways that
generate the chosen artifacts are entirely separate from the
upper-limb pathways that are damaged in these patients.}

Specifically:

\begin{itemize}[leftmargin=*]
    \item The eye blink is driven by the \textbf{facial nerve}
    (cranial nerve VII).
    \item Jaw clinching and bruxism are driven by the
    \textbf{trigeminal nerve} (cranial nerve V).
    \item Both are cranial nerves that exit the brainstem directly
    and never traverse the brachial plexus, the cervical/thoracic
    spinal segments serving the arm, or the peripheral nerves of the
    upper limb.
\end{itemize}

Consequently, an amputation at any level of the arm --- and any
associated damage to the brachial plexus or upper-limb peripheral
nerves --- leaves the blink and jaw pathways completely unaffected.
A patient who has lost all reliable myoelectric control of a residual
limb retains full, effortless voluntary control of blinking and jaw
movement. This dissociation is the core clinical argument for the
approach: it provides a high-bandwidth, reliable control channel that
is anatomically guaranteed to survive the very injuries that disable
conventional prosthetic interfaces.

This argument naturally has boundaries. Patients with concomitant
cranial nerve injury, certain facial palsies, severe bruxism-related
temporomandibular disorders, or neuromuscular conditions affecting the
face and jaw may not be suitable candidates. The approach is therefore
positioned as an option for the specific --- and clinically common ---
population of upper-limb amputees with intact craniofacial function,
rather than as a universal solution.

%% ----------------------------------------------------------
\section{Clinical Safety and Comfort Considerations}
\label{sec:med_safety}
%% ----------------------------------------------------------

Because the system requires the user to produce voluntary facial and
jaw gestures repeatedly during operation, the clinical safety and
ergonomic comfort of these gestures warrant explicit consideration.

\textbf{Eye blinks} are physiologically trivial: humans blink
spontaneously thousands of times per day, and voluntary blinking adds
negligible additional load. No safety concern arises from blink-based
commands.

\textbf{Jaw clinching and bruxism}, however, involve repeated forceful
contraction of the muscles of mastication. While occasional voluntary
clinching is harmless, sustained or excessive grinding is associated
in the clinical literature with temporomandibular joint (TMJ) strain,
masticatory muscle fatigue, and, in pathological (involuntary) bruxism,
dental wear. For a control interface, these considerations imply
several design safeguards:

\begin{itemize}[leftmargin=*]
    \item Commands should be brief and require only moderate force ---
    the activity gate (Chapter~7) is calibrated to detect deliberate
    gestures without demanding maximal contraction.
    \item The command vocabulary should minimise the frequency of
    jaw gestures relative to the lower-effort blink gestures where
    possible.
    \item Extended continuous use should incorporate rest periods to
    avoid masticatory muscle fatigue.
\end{itemize}

These are engineering-level mitigations of an inherently low clinical
risk for occasional, moderate use; they would require formal
evaluation before any sustained clinical deployment, which is outside
the scope of this proof-of-concept project.

%% ----------------------------------------------------------
\section{Clinical Rationale for Artifact-Based Control}
\label{sec:med_rationale}
%% ----------------------------------------------------------

Bringing the physiological arguments together, the case for an
artifact-based interface over motor imagery for this specific
population can be summarised along three dimensions:

\textbf{Signal reliability.} Motor-imagery rhythms are small
(1--\SI{5}{\micro\volt}), highly variable between individuals, and
absent altogether in a substantial fraction of users (the so-called
``BCI illiteracy'' phenomenon). Voluntary artifacts are large
(tens to hundreds of microvolts), reproducible within an individual,
and producible by essentially anyone with intact craniofacial
function. For a clinical assistive device that must work reliably for
its user every day, this difference is decisive.

\textbf{Calibration burden.} Motor-imagery systems typically require
lengthy, repeated calibration sessions and frequent recalibration as
the cortical signal drifts. Voluntary artifacts, being large and
stereotyped, require far less calibration data and are more stable
across sessions --- an important practical advantage for a device
intended for daily independent use.

\textbf{Anatomical robustness.} As established in
Section~\ref{sec:med_intact}, the cranial pathways generating these
artifacts are anatomically guaranteed to survive upper-limb
amputation and peripheral nerve injury. The interface therefore offers
a control channel whose availability is decoupled from the patient's
injury, which is precisely the property that myoelectric control
lacks for this population.

Taken together, these considerations make voluntary artifact-based
control not merely a convenient engineering choice but a clinically
motivated one for upper-limb amputees with compromised residual
musculature.

%% ----------------------------------------------------------
\section{Chapter Summary}
\label{sec:med_summary}
%% ----------------------------------------------------------

This chapter established the medical and physiological foundations of
the artifact-based control paradigm. Scalp electrodes capture a
mixture of cortical, ocular, and muscular signals; this project
deliberately exploits the ocular and muscular signals --- normally
discarded as artifacts --- as voluntary control commands. The eye
blink, driven by the facial nerve, produces a large corneoretinal
deflection at the frontal electrodes whose morphology encodes single
and double blink commands. Jaw clinching and bruxism, driven
by the trigeminal nerve, produce broadband \acrshort{emg} at the
temporal electrodes; the bilateral symmetry of the clinch and the
lateral asymmetry of bruxism allow the close, rotate-left, and
rotate-right commands to be distinguished, with the left/right
distinction grounded in the differential activation of the left and
right temporalis muscles. Critically, both the facial and trigeminal
pathways are cranial and therefore anatomically independent of the
upper-limb nerves damaged in the target population, making these
control signals robustly available even when myoelectric control is
not. With this physiological grounding established, the following
chapter describes the hardware that captures and acts upon these
signals.